\section{Обзор}
\label{sec:review}

Обзор охватывает четыре аспекта: семейство мер асимптотической относительной эффективности
(АОЭ) статистических критериев; аналитические результаты эффективности по
Ходжесу--Леману для критериев согласия; способы численного оценивания экспоненциальных
показателей; программные инструменты.
Задача обзора~--- показать место работы среди
существующих результатов и обосновать вычислительную схему, архитектуру и стек.

\subsection{Постановка задачи сравнения критериев согласия}
\label{ssec:statement}

Критерии согласия предназначены для проверки соответствия выборки заданному закону
распределения; во фреймворке \texttt{PySATL} реализованы десятки таких критериев различных
типов, что делает актуальной задачу их объективного сопоставления. Сравнение при конечном
объёме выборки зависит от уровня значимости и от самого объёма, поэтому в теории
используются асимптотические меры эффективности, каждая из которых относится к своему
режиму альтернатив~[7].

Под эффективностью в данной работе понимается способность критерия при увеличении объёма
выборки распознавать заданную альтернативу при фиксированном уровне значимости. Для
обоснованного выбора меры необходимо сопоставить существующие подходы и способы их
вычисления, применявшиеся в теории и на практике, а также определить, какие из них
реализуемы для критериев, уже присутствующих в библиотеке.

\subsection{Меры асимптотической относительной эффективности}
\label{ssec:measures}

\textbf{Эффективность Питмена} рассматривает локальные альтернативы, приближающиеся к
нулевой гипотезе со скоростью $1/\sqrt{n}$, и сравнивает объёмы выборок, необходимые для
достижения одинаковой мощности. Мера опирается на асимптотическую нормальность статистики и
потому хорошо описывает поведение критериев в окрестности $H_0$, однако теряет
содержательность для статистик с невырожденными неасимптотически нормальными пределами,
к которым относится большинство критериев согласия, и не характеризует критерии на
фиксированных, содержательно значимых альтернативах~[7].

\textbf{Эффективность Бахадура} фиксирует альтернативу и сравнивает критерии по скорости
убывания достигнутого уровня значимости, то есть ошибки первого рода; точный наклон
выражается через функцию больших уклонений статистики при нулевой гипотезе~[8]. Мера
получила широкое распространение, однако для статистик с невырожденным предельным
распределением точный наклон вырождается, и применяется приближённый~[9]; его вывод
нетривиален и, как правило, требует индивидуального теоретического анализа для каждого
критерия. В смежной работе эту меру реализует коллега по команде.

\textbf{Эффективность по Ходжесу--Леману}~[1] фиксирует альтернативу $H_1(\theta)$, ограничивает уровень значимости
$\alpha$ и сравнивает критерии по экспоненциальной скорости убывания вероятности ошибки
второго рода. Если для критерия $T_n$ при больших $n$ выполнено соотношение
\begin{equation}
  1 - \pi_n(\theta) \;\approx\; C\,\exp\!\bigl(-n\,\rho(\theta)\bigr),
  \qquad n \to \infty,
  \label{eq:hl-decay}
\end{equation}
где $\pi_n(\theta)$~--- мощность критерия, то величина
\begin{equation}
  \rho(\theta) \;=\; -\lim_{n\to\infty}\frac{1}{n}\,
  \ln\!\bigl(1-\pi_n(\theta)\bigr) \;\ge\; 0
  \label{eq:hl-rho}
\end{equation}
называется \emph{показателем скорости} (показателем эффективности по Ходжесу--Леману), а
относительная эффективность двух критериев $K_1$, $K_2$ определяется как
$\operatorname{ARE}_{K_1/K_2}(\theta) = \rho_1(\theta)/\rho_2(\theta)$. Показатель
$\rho(\theta)$ определяется большими уклонениями статистики \emph{при альтернативе}, а не
при нулевой гипотезе, что и отличает данную меру от бахадуровской. В первоисточнике~[1]
отмечено, что мощность не обязана быть выпуклой функцией объёма выборки, поэтому
классическое определение относительной эффективности через отношение необходимых объёмов
выборок требует оговорок и сравнивать следует именно показатели скорости. Общая теория изложена
в работах Чернова и Калленберга~[2, 3], современная формулировка для общих тестовых
статистик~--- в~[4]; меры Бахадура и Ходжеса--Лемана характеризуют разные события, поэтому
их значения в общем случае не совпадают, и сопоставление двух реализаций на общих данных
служит независимой проверкой согласованности~[8].

\subsection{Вычисление показателя скорости}
\label{ssec:routes}

Теоретическое исследование эффективности по Ходжесу--Леману опирается на аппарат теории
больших уклонений. Ключевую роль сыграли результаты Чернова~[2], связавшие скорость
экспоненциального убывания ошибки с вариационным функционалом; в дальнейшем подход
обобщался и уточнялся~[3, 4, 5].

Для конкретных классов критериев согласия аналитические выражения показателя скорости
получены, но точечно: в монографии Никитина вычислены показатели для критериев согласия
основных параметрических семейств, включая замкнутые выражения для статистик интегрального
и экстремального типа~[5], в работе Янссена построены глобальные функции мощности
критериев согласия и показана оптимальность критерия Колмогорова--Смирнова по
Ходжесу--Леману~[6]. Для составной нулевой гипотезы с оцениваемыми параметрами ситуация
хуже: предельные распределения статистик эмпирической функции распределения описаны в
терминах собственных значений интегральных операторов~[9], однако большие уклонения этих
статистик при альтернативе в общем виде не опубликованы. Для большинства критериев,
реализованных в прикладных библиотеках, явные формулы либо трудоёмки для вывода, либо
отсутствуют, и само их вычисление требует аккуратного обращения с предельными
распределениями, условиями регулярности и асимптотикой хвостов.

Прямой численный метод состоит в следующем: ошибка второго рода
$\beta_n(\theta) = 1-\pi_n(\theta)$ оценивается моделированием на сетке объёмов $n$,
зависимость $\ln\hat\beta_n$ от $n$ аппроксимируется прямой, и угловой коэффициент даёт
оценку $\hat\rho(\theta)$. Ограничение такого метода~--- разрешимость: при $M$ итерациях оценка
отличима от нуля лишь пока $\beta_n \gtrsim 10/M$, поэтому сетка объёмов образует окно,
зависящее от $M$, $\alpha$ и $\theta$; погрешность оценки контролируется биномиальной
дисперсией и стандартными процедурами выбора числа повторов~[13]. Корректность оценки
мощности при фиксированном уровне значимости дополнительно требует аккуратного вычисления
порога отвержения по предельному либо смоделированному при нулевой гипотезе распределению,
поскольку ошибка второго рода при больших объёмах выборки является редким событием~[12].
В публикациях по эффективности критериев моделирование используется вспомогательно~--- для
критических значений и мощностей~[9, 10], но не для извлечения показателя скорости.

\subsection{Программные инструменты и реализации}
\label{ssec:tools}

Анализ доступных инструментов показывает следующее. В пакете \texttt{R} встречаются
отдельные решения для сравнения тестов (например, связанные с бутстрепом, ранговыми процедурами),
однако, как правило, оценка асимптотической относительной эффективности реализуется в виде
разовых скриптов, ориентированных на конкретную пару критериев и конкретную альтернативу.

В экосистеме \texttt{Python} инструменты типа \texttt{statsmodels}, \texttt{scipy.stats},
\texttt{pingouin} предоставляют средства для расчёта мощности критерия, но не предлагают
унифицированного механизма для \emph{эксперимента} по оценке эффективности по Ходжесу--Леману
с автоматизированным построением оценок $\rho(\theta)$, матрицы относительных эффективностей и
последующим оформлением результатов.

\subsection{Новизна работы}
\label{ssec:novelty}

Из проведённого обзора следуют четыре аспекта новизны.

\begin{itemize}
  \item \textbf{Формализация в виде типа эксперимента.} Оценка эффективности критериев
        согласия по Ходжесу--Леману реализуется как самостоятельный тип эксперимента
        фреймворка \texttt{pysatl-experiment}, а не как разрозненные скрипты.
  \item \textbf{Численный метод извлечения показателя скорости.} Показатель $\rho(\theta)$
        оценивается по лог-линейной аппроксимации ошибки второго рода на окне разрешимости,
        что в источниках не встречается: моделирование используется в них вспомогательно,
        для критических значений и мощностей~[9, 10].
  \item \textbf{Подменяемый этап вычисления.} Способ вычисления $\rho(\theta)$~---
        аналитический либо численный~--- выбирается в конфигурации эксперимента, что
        позволяет использовать готовые теоретические результаты при их наличии и наращивать
        аналитические возможности поэтапно.
  \item \textbf{Интеграция в существующий конвейер.} Эксперимент встраивается в общую
        логику \texttt{pysatl-experiment}, что обеспечивает
        сопоставимость с имеющимися типами экспериментов.
\end{itemize}

\subsection{Выводы}
\label{ssec:conclusions}

\begin{enumerate}
  \item \textbf{Вычислительная схема.} Основной способ~--- прямое моделирование ошибки
        второго рода на окне разрешимости с лог-линейной регрессией и доверительными
        интервалами~[13]; аналитический способ сохраняется как подменяемая стратегия.
  \item \textbf{Валидация.} Эталоном служат аналитические значения~[5, 6]; при их
        отсутствии~--- контрольные случаи с известным показателем: для простой нулевой
        гипотезы тест отношения правдоподобия воспроизводит расстояние Кульбака--Лейблера
        как следствие экспоненциальных оценок Чернова~[2].
  \item \textbf{Архитектура.} Новый тип эксперимента \texttt{hodges\_lehmann\_efficiency}
        строится над существующими генерацией выборок и кэшем предельных распределений;
        предусматривается архитектурное решение, позволяющее выбирать способ вычисления
        показателя скорости: либо на основе имеющихся аналитических соотношений, либо
        численным путём; результаты агрегируются в оценки $\rho(\theta)$ и матрицу $\operatorname{ARE}$.
  \item \textbf{Цель и задачи.} Результаты обзора подтвердили корректность ранее
        сформулированных цели и задач.
\end{enumerate}

\clearpage